\documentclass{ctexart}

% === 页面布局 ===
\usepackage[left=2.5cm, right=2.5cm, top=2.5cm, bottom=2.5cm]{geometry}

% === 数学与符号 ===
\usepackage{amsmath, amssymb}

% === 链接 ===
\usepackage[colorlinks,linkcolor=black,citecolor=black,urlcolor=black]{hyperref}

% === 定理类环境 ===
\newtheorem{theorem}{定理}[section]
\newtheorem{lemma}[theorem]{引理}
\newtheorem{corollary}[theorem]{推论}
\newtheorem{definition}[theorem]{定义}
\newtheorem{proof}[theorem]{证明}

% === 标题 ===
\title{\bfseries 经典NP定义的语义缺失与操作不自足：\\
	为什么P vs.\ NP问题需要一个新的概念框架}
\author{郑波尽\thanks{中南民族大学计算机学院，武汉，430074，中国。Email: \texttt{zhengbojin@scuec.edu.cn}。}
	\and 郑景文\thanks{武汉大学网络安全学院，武汉，430072，中国。}
	\and 王伟武\thanks{国家开发银行信息科技处，北京，100033，中国。}
	\and 黄芷莹\thanks{中南民族大学计算机学院，武汉，430074，中国。}
}
\date{2026年7月}

\begin{document}
	
	\maketitle
	
	\begin{abstract}
		$\mathbf{P}$ versus $\mathbf{NP}$问题困扰理论计算机科学已逾半个世纪。相对化、自然证明和代数化三大障碍的存在表明，问题的困难可能不仅来自计算本身，更来自我们用以描述计算的数学语言。本文提出一个哲学诊断：经典NP定义存在根本性的语义缺失与操作不自足。该定义将非确定性的本质外包给外部量词``存在见证''$\exists c$和外部谓词$IsWitness(c)$，而经典确定性图灵机的内部语义无法内在地表达该外部谓词所必需的不可公度性（即生成元在基域上的线性无关性）。更为根本的是，经典NP定义隐式地要求了一个在经典操作语义中\textbf{不可实现}的操作——将虚部控制信息拷贝到实部成为数据$c$。本文论证了这一拷贝操作在经典图灵机模型中的不可能性，并指出经典NP定义因此是\textbf{操作不自足}的：它诉诸了一个其自身形式体系无法提供的资源。在此基础上，我们进一步诊断了经典NP定义的三重概念错误，它们正是这一拷贝不可能性的必然结构性后果——第一重错误（函数指针被强制转换为数据）导致\textbf{不可公度性的丢失}：见证$c$的各个比特在虚部中对应于不同的不可公度生成元选择，转换为数据比特后，比特间不再保留任何代数独立性关系；第二重错误（编码为比特串并用存在量词包裹）导致\textbf{对偶性的丢失}：$+\sqrt{p_i}$与$-\sqrt{p_i}$的伽罗瓦对偶关系在比特$0$和$1$中没有对应物；因为不可公度性和对偶性都已丢失，经典模型无法正面刻画非确定性选择的独立性和结构，只能将其外包给外部量词和外部谓词，导致非确定性概念的\textbf{语义残缺}（第三重错误）。基于这些诊断，我们分析了``聪明算法''假设中隐含的语义错位：它要求用无法表达不可公度性且丧失了对偶性的内部语言，去等价地实现依赖于不可公度性和对偶性的外部语义——更精确地说，它要求用\textbf{无法执行虚部到实部拷贝的模型}，去等价地实现一个其定义本身\textbf{依赖于这一拷贝结果}的验证过程。这一诊断与哥德尔不完备定理和塔斯基真不可定义定理有着深刻的家族相似性，共同指向形式系统的自指性局限。最后，我们提出``数计一体''的方法论原则，并概述了CBTM/IVM新框架如何通过参数化设计——将不可公度性内建为操作语义的内在属性、让对偶性常驻虚部、从根本上消除拷贝操作的必要性——为消解上述困境提供了概念基础。
		
		\vspace{0.5cm}
		\noindent \textbf{关键词：} $\mathbf{P}$ versus $\mathbf{NP}$，语义缺失，操作不自足，不可公度性，对偶性，虚部到实部拷贝，操作语义，数计一体，哥德尔不完备定理
	\end{abstract}
	
	\section{引言}
	
	\subsection{P vs.\ NP问题：半个世纪的困境}
	
	$\mathbf{P}$ versus $\mathbf{NP}$问题\cite{cook1971,levin1973}是理论计算机科学乃至整个现代数学中最具影响力的未解难题之一。$\mathbf{P} \subseteq \mathbf{NP}$是平凡的，但$\mathbf{P}$是否等于$\mathbf{NP}$——即所有可以在多项式时间内被验证的问题是否也可以在多项式时间内被求解——至今未解。
	
	经过半个多世纪的努力，研究者们逐渐识别出三个经典的证明障碍。Baker、Gill和Solovay~\cite{baker1975}的相对化障碍证明，存在两个谕示世界，一个使$\mathbf{P}=\mathbf{NP}$，另一个使$\mathbf{P}\neq\mathbf{NP}$，任何对谕示``透明''的证明技术都无法同时兼容这两个世界。Razborov和Rudich~\cite{razborov1997}的自然证明障碍指出，任何同时满足``有用性''、``大尺度''和``构造性''的组合下界证明将导致伪随机数生成器的攻破，因而被认为几乎不可能。Aaronson和Wigderson~\cite{aaronson2009}的代数化障碍将相对化推广为代数化，证明任何对谕示多项式扩展保持有效的证明技术同样无法分离$\mathbf{P}$和$\mathbf{NP}$。
	
	这三大障碍共同揭示了一个深层问题：传统计算模型是``语义贫乏''的——它们只关心符号的语法操作，抛弃了符号背后的数学语义。语义的缺失使得任何纯语法的论证都容易被外部谕示或伪随机函数所愚弄。障碍的存在暗示，问题的难度可能不仅来自计算本身，更来自我们用以描述计算的数学语言。
	
	\subsection{超越技术障碍：概念适当性的问题}
	
	在思考P vs.\ NP问题的历史中，有一个关键区分往往被忽视：\textbf{技术困难}与\textbf{概念困难}之间的区别。
	
	技术困难是指在给定框架内寻找证明的困难。例如，在欧几里得几何中证明平行公设是一个技术困难——数学家们花费了两千多年试图从其他四条公设推导出第五公设。概念困难则是指框架本身是否适合提出该问题的困难。平行公设问题最终被证明不是一个技术困难，而是一个概念困难：欧几里得几何的语言不足以表达非欧几何的可能性，需要重新定义``平行''的概念。
	
	本文试图论证，P vs.\ NP问题同样面临一个概念困难。具体而言，经典NP定义存在一个根本性的语义缺失与操作不自足，使得该问题在经典框架内难以获得充分的表达资源。这不是说该问题没有真值，也不是说我们从$\mathbf{P}=\mathbf{NP}$可以推导出一个形式逻辑矛盾。而是说，\textbf{经典计算模型的语言缺乏正面定义非确定性选择独立性所需的核心概念资源，且经典NP定义要求了一个在该模型操作语义中不可实现的操作}，因此任何局限于经典语言内的证明都难以触及P与NP差异的本质。
	
	本文是一篇哲学诊断论文，而非数学证明。本文的任务是阐明为什么经典框架无法完成这一证明，以及为什么新框架的构建是必要的。
	
	\subsection{论证结构概述}
	
	本文的论证分为三个层次。
	
	\textbf{第一层次（第2-5节）：诊断经典NP定义的语义缺失与操作不自足。} 我们首先阐明计算中语法与语义的分离，并区分两种意义上的``知道''一个数学属性。随后，我们分析经典NP定义的逻辑结构，指出它将非确定性的本质进行了``双重外包''：外包给外部量词``存在见证''和外部谓词$IsWitness(c)$。在此基础上，我们论证了一个被长期忽视的关键事实：经典NP定义隐式地要求了一个在经典操作语义中不可实现的``虚部到实部拷贝''操作，这使得经典NP定义是操作不自足的。我们进一步诊断了由此必然导致的三重概念错误：第一重错误（函数指针被强制转换为数据）导致不可公度性的丢失；第二重错误（编码和量词）导致对偶性的丢失；因两者皆已丢失，经典模型只能将非确定性的本质外包，导致语义残缺（第三重错误）。我们论证，外部谓词$IsWitness(c)$的内涵需要不可公度性来填充，而不可公度性恰恰是经典模型无法在强意义上表达的。
	
	\textbf{第二层次（第6-7节）：确立经典模型的表达边界。} 我们给出不可公度性元定理的完整证明，并从语法不变性和代数封闭性两个互补视角加以论证。基于这一元定理以及拷贝操作的不可能性，我们分析``聪明算法''假设中隐含的语义错位。我们将这一诊断与哥德尔不完备定理和塔斯基真不可定义定理进行比较，揭示两者共享的深层结构。
	
	\textbf{第三层次（第8-10节）：从诊断到解决。} 我们比较BSS模型与CBTM/IVM框架，阐明``内建语义''与``扩展域''两条路径的本质区别。我们提出``数计一体''的方法论原则，并概述CBTM/IVM新框架如何通过参数化设计来消解上述困境：将不可公度性内建为操作语义的内在属性，让对偶性常驻虚部，从根本上消除拷贝操作的必要性，使得经典NP定义中的三重概念错误和操作不自足得到系统性修正。最后，我们总结诊断，并阐明从``消解''到``解决''的逻辑链条。
	
	\section{计算中的语法与语义}
	
	\subsection{图灵机的语法中立性}
	
	经典图灵机模型~\cite{turing1936}的核心特征是彻底的\textbf{语法中立性}。一台图灵机是一个七元组$M = (Q, \Sigma, \Gamma, \delta, q_0, q_{\text{accept}}, q_{\text{reject}})$，其中转移函数$\delta: Q \times \Gamma \to Q \times \Gamma \times \{L, R\}$是整个模型的灵魂。转移函数的核心特征是：其输入中的符号参数仅仅是$\Gamma$中的一个\textbf{标识}。$\delta$无法区分两个具有不同外部语义但形式上可互换的符号，它只能根据当前状态和当前符号的标识来决定下一步动作。
	
	同一个二进制串``101''可以同时被解释为数字5、布尔序列``真、假、真''或某个代数扩张中的元素$a+b\sqrt{2}$的编码。计算的过程是符号操作，但计算的``意义''完全存在于外部观察者的解释之中。
	
	正是这种语法中立性赋予了图灵机通用模拟的能力——它可以模拟任何算法，因为它不预设任何特定的语义。但也正是这种语法中立性，埋下了一个深层的隐患：当我们需要区分计算的``难度''时，纯粹语法层面的描述可能丢失了决定难度的关键数学信息。
	
	\subsection{数学计算作为对比：语法与语义的统一}
	
	数学计算天然携带丰富的语义信息。当我们讨论一个代数数$a + b\sqrt{2}$时，我们关心的不仅是符号书写的形式，更是它在扩域$\mathbb{Q}(\sqrt{2})$中的代数性质：它与有理数的不可公度性、它作为二维向量空间中的元素的位置，等等。
	
	不可公度性是一个深刻的数学概念。在数学中，两个代数数$\gamma, \delta$在有理数域$\mathbb{Q}$上\textbf{不可公度}，当且仅当它们在$\mathbb{Q}$上线性无关，即不存在不全为零的有理数$c_1, c_2$使得$c_1\gamma + c_2\delta = 0$。这意味着$\gamma$和$\delta$之间不存在任何有理线性关系，它们各自贡献了一个独立于$\mathbb{Q}$的代数维度。
	
	数学计算的每一步都可以被赋予语义解释。当我们写下$(a + b\sqrt{2}) \cdot (c + d\sqrt{2})$并展开时，我们不仅在操作符号，同时也在执行$\mathbb{Q}(\sqrt{2})$中的域乘法。语法操作与语义解释始终保持着严格的对应关系，数学推理的正确性正是依赖于这种一致性。
	
	\subsection{两种意义上的``知道''一个数学属性}
	
	我们区分两种意义上的``知道''一个数学属性。
	
	\subsubsection{弱意义：算法可判定性}
	
	在弱意义上，我们说一台机器``知道''一个数学属性，如果它在给定适当编码的输入后，能够输出关于该属性的正确判定。
	
	例如，经典图灵机可以在适当的编码下判定两个代数数是否在$\mathbb{Q}$上线性无关。给定两个代数数$\gamma, \delta$的适当表示（例如，用它们的极小多项式和隔离区间），存在多项式时间算法（如高斯消元）能够判定它们是否线性无关。不可公度性本身在$\mathbf{P}$内是可计算的——这是一个重要的数学事实。
	
	然而，这种``知道''依赖于外部编码和语义赋值。编码函数$\varepsilon$将数学对象映射为符号串，机器操作这些符号串，结果再由外部观察者解码回数学命题。机器本身并不``理解''它所操作的对象具有不可公度性——它只是在执行语法规则。
	
	\subsubsection{强意义：操作语义的内蕴性}
	
	在强意义上，我们说一台机器``知道''一个数学属性，如果该属性是机器操作语义的\textbf{内在属性}。一个属性$\mathfrak{P}$是图灵机操作语义的内在属性，当且仅当它完全由机器的转移规则和符号标识决定，而不依赖于外部观察者对符号的语义赋值。
	
	形式化地，$\mathfrak{P}$是内在属性当且仅当对于字母表上任意保持空白符不变的置换$\pi$，$M$在输入$w$上具有$\mathfrak{P}$当且仅当$M_\pi$在输入$\pi(w)$上具有$\mathfrak{P}$。例如，``机器在输入$w$上停机''是一个内在属性——无论我们将符号解释成什么数学对象，停机的判断只取决于操作规则本身。
	
	\subsubsection{关键区分：判定不可公度性与内在地锚定不可公度性}
	
	\textbf{本文的核心论断之一是：经典图灵机可以在弱意义上判定不可公度性，但无法在强意义上内在地锚定它。}
	
	不可公度性元定理（将在第6节中给出完整证明）并不否认经典图灵机可以通过适当的编码来判定两个代数数是否可公度。它证明的是一个更强的结论：\textbf{图灵机的操作语义本身无法内在地``感知''它所处理的符号是否具有不可公度性。}即，机器的操作行为不能唯一地确定符号的外部解释必须包含不可公度元，因为符号可以被置换，而机器行为在置换下保持不变。这一区分的理论分量在于：P vs NP问题的困难不在于``不可公度性是否可判定''，而在于``经典模型的操作语义是否能够正面定义非确定性选择的独立性''。后者恰恰需要不可公度性作为操作语义的内在属性，而非仅仅作为外部编码的语义赋值。
	
	\subsubsection{非平凡编码的影响}
	
	在弱意义下，经典图灵机需要通过非平凡编码来判定不可公度性。但非平凡编码本身可能改变计算能力或计算复杂性。
	
	考虑多带图灵机可被单带图灵机多项式时间模拟的经典定理。这个定理本身是正确的，它是一个纯语法层面的结论。然而，当我们将它应用于具有代数语义的数学对象时，编码过程可能产生非中立的后果。复合符号$(a, b)$的数学语义是什么？它可以是有理数$a/b$，也可以是代数数$a + b\alpha$，还可以是实数的一个近似表示。图灵机本身无法区分这些语义，而模拟过程也从未保证这些语义的一致性。
	
	两个思想实验可以具体展示这种语义丢失。实验一：连续量的离散化。实数是连续的，任何有限编码都只能表示其近似值，模拟器操作的只是一个实数的``影子''，而非实数本身。实验二：不可公度性的抹除。$\sqrt{2}$在有理数域上不可公度，这一代数事实在编码为有序对$(a, b)$后完全消失。单带机看到的只是两个有理系数，它无从知晓$b$代表的是$b\sqrt{2}$而非有理数$b$。
	
	这两个实验共同揭示了一个深刻的洞见：\textbf{语法模拟的成功并不意味着语义的保持}。编码的本质是外延指代，而非内涵内建。它将不可公度性``表示''为符号串，但未将其``内建''为操作语义。
	
	\section{经典NP定义的双重外包}
	
	\subsection{P的内在定义}
	
	在经典计算理论中，$\mathbf{P}$的定义是内在的、操作性的。一个语言$L$属于$\mathbf{P}$，当且仅当存在一台多项式时间确定性图灵机$M$，使得对于任意输入$x$，$x \in L$当且仅当$M$在输入$x$上停机于接受状态。
	
	这个定义的关键特征是：$\mathbf{P}$的语义完全由机器的操作语义提供。``多项式时间内可判定''是一个内在属性——它完全由转移函数、状态集和纸带操作规则决定，不依赖于任何外部语义赋值。
	
	\subsection{NP的外在定义}
	
	与$\mathbf{P}$形成鲜明对比的是，$\mathbf{NP}$的定义是外在的、外延性的：
	
	\[
	x \in L \iff \exists c \; [IsWitness_L(c) \land (|c| \le p(|x|) \land V(x, c) = 1)].
	\]
	
	这一定义中，$IsWitness_L(c)$是一个独立的谓词，断言``$c$是关于语言$L$的一个见证''，即$c$具有作为见证的资格或本质。验证器$V$本身是确定性的——它在输入$(x, c)$上的行为完全由转移函数决定。非确定性被完全外包给了这一定义中的两个外部元素。
	
	我们将这一外包结构称为\textbf{双重外包}。
	
	\textbf{第一重外包：外部量词``$\exists c$''。} 非确定性选择的存在性被外包给了元语言的存在量词。机器本身不``知道''什么是``存在''——这个量词属于我们用来讨论图灵机的数学语言（元语言），而非机器内部的操作语言。验证器$V$只是确定性地检查一对$(x, c)$是否满足接受条件，但$c$的``存在''从未在机器内部被构造或定义。
	
	\textbf{第二重外包：外部谓词$IsWitness_L(c)$。} ``什么是见证''这一内涵属性被外包给了外部谓词。在经典框架中，这个谓词的语义被外延性地还原为``存在验证器$V$使得$V(x,c)=1$''——即$c$是见证，当且仅当验证器接受。但``是见证''这一内在属性——即$c$与$x$之间的内在代数关系，为什么$c$构成了$x$的一个解——从未被正面定义。
	
	\subsection{为什么外部谓词需要不可公度性}
	
	$IsWitness_L(c)$需要回答一个根本问题：\textbf{为什么$c$是$x$的见证？}这要求正面刻画见证与输入之间的内在代数关系，即非确定性分支选项之间的独立性。
	
	在CBTM（复布尔图灵机）框架中，这一独立性被精确刻画为不可公度性。见证$c$本质上是一个函数指针的地址序列。每一个比特对应于一次非确定性选择——从两个互补的函数指针中选取一个。这些函数指针对应于扩域$\mathbb{F}_4/\mathbb{F}_2$中的生成元$\alpha$和$\beta$，它们在基域$\mathbb{F}_2$上不可公度。正是这种代数上的不可公度性，保证了不同选择路径之间的独立性——它们无法通过基域内的任何运算相互转化。
	
	因此，$IsWitness_L(c)$的内涵需要不可公度性来填充。见证之所以是见证，是因为它的每一个比特都代表了一个不可公度的选择——一个无法被确定性操作所模拟或消除的、真正的非确定性分支。
	
	\subsection{经典定义中缺失的前提}
	
	经典NP定义以一个未经审视的预设为前提：\textbf{经典模型的内部语义足以正面刻画非确定性选择的独立性。}这个预设从未被正面检验。经典教科书将``有多个后继''视为非确定性的全部，将选项之间的``独立性''视为不言自明的直觉前提。但为什么这些选项是``独立的''？它们的独立性由什么来保证？标准定义对此保持沉默。
	
	本文的诊断是：这个预设是假的。经典图灵机的操作语义无法内在地表达不可公度性，而非确定性选择的独立性恰恰需要不可公度性来保证。因此，经典NP定义在语义上是不完整的：它指明了哪些问题是NP的（外延），但没有正面刻画非确定性的本质是什么（内涵）。外部谓词$IsWitness_L(c)$是一个语义空洞——它指向了一个经典语言无法表达的核心概念。
	
	这一语义缺失正是三大障碍能够长期存在的根本原因。相对化障碍的根源在于：谕示可以被暗中赋予填补语义空洞的超能力——替机器完成“猜测虚部控制信息并拷贝到实部”这一经典模型自身不可能执行的操作。从CBTM架构来看，虚部与实部分属控制层与数据层，经典机器无法感知虚部，更无法将虚部信息跨层次拷贝到实部；但谕示可以免费提供这一拷贝服务。当这一拷贝操作由谕示免费完成，P与NP之间的鸿沟便在谕示加持下消失了——这正是相对化世界中出现$\mathbf{P}^A = \mathbf{NP}^A$的根源。自然证明障碍利用的是自然性质基于组合属性，而独立性的本质是代数概念，超出了组合范畴。代数化障碍利用的是代数化谕示只能回答多项式等式问题，无法感知符号的正负号——而正负号正是对偶性在操作层面的直接体现，对偶性在经典NP定义中已被存在量词消除。
	
	\section{拷贝操作的不可能性：经典NP定义的操作不自足}
	
	第3节的诊断揭示了经典NP定义将非确定性的本质外包给外部量词和外部谓词。在进入三重概念错误的详细分析之前，我们需要首先确立一个更为根本的事实：\textbf{经典NP定义隐式地要求了一个在经典图灵机操作语义中不可实现的操作。}
	
	\subsection{见证$c$的本体论来源：虚部控制信息}
	
	在CBTM/IVM框架的视角下，见证$c$的本质是\textbf{虚部上的控制信息}。非确定性验证过程的语义表明，$c$的每一个比特实质上指示验证器在每一次非确定性选择点上是走激活路径（选择$+\sqrt{p_i}$）还是不激活路径（选择$-\sqrt{p_i}$）。$c$因此是一个\textbf{函数指针的地址序列}——它规定了验证器在计算树中的行走路径，告诉机器``执行哪条指令''，而非``操作什么数据''。
	
	然而，经典NP定义将$c$写在纸带上，作为\textbf{实部数据}供验证器$V$读取。这意味着：
	
	\begin{quote}
		\textbf{经典NP定义隐式地假设了一个从虚部到实部的``拷贝''操作：将控制信息（函数指针地址序列）转换为数据（纸带上的比特串）。}
	\end{quote}
	
	\subsection{投影约束公理与拷贝操作的不可实现性}
	
	CBTM/IVM框架中的投影约束公理\cite{cbtmarxiv}规定：
	\begin{enumerate}
		\item 写入符号的虚部完全由当前状态和读取符号的虚部通过布尔函数决定。机器不能``凭空''产生虚部非零的符号，也不能将实部信息``提升''为虚部信息。
		\item 特别地，若读取符号的虚部为零，则写入符号的虚部也必须为零。这意味着，一旦计算处于基域（纯实部），它就无法产生涉及不可公度元的符号。
	\end{enumerate}
	
	在经典图灵机模型中，纸带符号是纯语法的——它们没有虚部/实部的区分。但从CBTM的视角回看，经典图灵机只能操作``虚部恒为零''的符号。经典模型的所有操作都停留在基域内，它无法：
	\begin{itemize}
		\item \textbf{感知}虚部的存在（不可公度性元定理，见第6节）；
		\item \textbf{生成}虚部非零的符号（因为没有域扩张机制）；
		\item \textbf{将虚部信息拷贝到实部}（因为这是跨语义层次的操作——将控制信息转换为数据，需要首先感知控制信息的存在）。
	\end{itemize}
	
	因此，经典NP定义中``将虚部控制信息拷贝到实部成为数据$c$''这一操作，在经典图灵机的操作语义中是\textbf{不可实现的}。
	
	\subsection{操作不自足的含义}
	
	这一不可能性揭示了经典NP定义的根本问题：它是\textbf{操作不自足}的——它诉诸了一个其自身形式体系无法提供的资源。
	
	在实际的经典NP验证中，$c$从不是由机器内部产生的。$c$是由\textbf{外部观察者（元语言）直接提供的}。外部观察者``猜测''了非确定性选择的轨迹，将其编码为比特串，写在纸带上，然后让验证器读取。这意味着：经典NP定义中的``非确定性''根本不是机器内部的操作特征，而是\textbf{外部观察者赋予的语义标签}。量词``$\exists c$''所做的工作——猜测正确的见证——完全发生在机器外部，在元语言层面。机器本身只是一个确定性的验证器，它不知道什么是``存在''，什么是``见证''。
	
	CBTM/IVM框架将这一被掩盖的事实显式化：在CBTM中，非确定性分支由虚部$b=1$触发，分支选择（激活/不激活）由机器内部的操作语义执行，生成元的激活由分支-激活引理强制性记账。非确定性成为机器内部的操作特征，而非外部观察者的语义标签。\textbf{从虚部到实部的拷贝这一不可能操作被彻底消除。}
	
	\subsection{拷贝不可能性与三重概念错误的关系}
	
	第5节将详细诊断经典NP定义的三重概念错误。这里先行指出，这些概念错误不是偶然的疏忽，而是\textbf{拷贝操作不可能性的必然结构性后果}：
	
	\begin{enumerate}
		\item 因为拷贝不可实现，经典NP定义只能将$c$直接写在实部纸带上，从而将函数指针\textbf{强制扭曲}为数据。这一转换同时导致\textbf{不可公度性的丢失}：见证$c$的各个比特在虚部中对应于不同的不可公度生成元选择（如第$i$个比特对应选择$\sqrt{p_i}$还是$-\sqrt{p_i}$），这些生成元在基域上线性无关，保证了不同选择路径在代数上不可相互归约。一旦被转换为数据比特$0$和$1$，比特之间就不再保留任何代数独立性关系——$0$和$1$在基域内完全可以相互转化。\textbf{非确定性选择的独立性因此在概念层面被抹除}（第一重错误）。
		
		\item 从虚部到实部的编码\textbf{必然丢失对偶性}：在虚部中，每次选择是在$+\sqrt{p_i}$与$-\sqrt{p_i}$之间做出的，这两个对偶生成元由伽罗瓦自同构精确映射。一旦被编码为比特$0$和$1$，这种对偶关系就完全消失了——$0$和$1$在语法上没有对偶关系。存在量词$\exists c$进一步抹去了``未被选择的另一条路径''的所有痕迹（第二重错误）。
		
		\item \textbf{因为不可公度性已在第一重错误中丢失，对偶性已在第二重错误中丢失}，经典模型的操作语义无法正面刻画非确定性选择的独立性（需要不可公度性）和结构（需要对偶性），因此只能将非确定性的语义本质\textbf{外包}给外部量词和外部谓词。\textbf{第三重错误是前两重错误在语义层面的必然后果}。
	\end{enumerate}
	
	因此，经典NP定义的问题不是``做得不够好''，而是\textbf{在其自身框架内原则上不可自我完善}——任何局限于经典操作语义的NP定义都必然面临这些结构性缺陷。
	
	\section{经典NP定义的三重概念错误}
	
	第3节的诊断揭示了经典NP定义将非确定性的本质外包给外部量词和外部谓词。第4节确立了这一外包过程中隐含的拷贝操作在经典操作语义中的不可实现性。本节将详细展开由此必然导致的三重概念错误。这三重错误不是偶然的疏忽，而是拷贝不可能性——即经典模型无法将虚部控制信息拷贝到实部——在概念层面的必然结构性映射。
	
	\subsection{第一重错误：函数指针被强制转换为数据，导致不可公度性丢失}
	
	在经典NP定义中，见证$c$被理解为一个字符串，写在验证器的纸带上，供其读取和处理。按照经典图灵机的架构，纸带上的内容属于\textbf{数据范畴}——它们是被操作的对象。因此，$c$在经典框架中被默认为数据。
	
	然而，$c$的本质是\textbf{控制流信息}。非确定性验证过程的语义表明，$c$的每一个比特实质上指示验证器在每一次非确定性选择点上应当选取哪一条分支。换言之，$c$扮演的是\textbf{函数指针地址序列}的角色：它告诉机器``执行哪条指令''，而非``操作什么数据''。
	
	在计算系统的基本架构中，函数指针和数据是两种截然不同的对象：
	\begin{itemize}
		\item \textbf{函数指针（代码）}：决定控制流，告诉机器在多个可能的后续操作中选择哪一个。
		\item \textbf{数据}：被操作的对象，被机器读写和运算。
	\end{itemize}
	
	经典图灵机架构的一个基本原则是\textbf{代码与数据的严格分离}。转移函数（代码）是固定的，定义了机器的全部行为能力；纸带内容（数据）是可读写的，是被操作的对象。控制流由转移函数根据当前状态和读取符号唯一决定，而非由纸带上的数据直接控制。
	
	见证$c$恰恰打破了这一分离。$c$是纸带上的一个字符串——按照经典架构，它应该属于数据范畴。然而，$c$承担的却是\textbf{决定控制流}的代码功能——它规定了验证器在计算树中的行走路径，决定了哪些状态序列被激活、哪些函数指针被调用。经典NP定义将这一函数指针地址序列强制写入纸带，完成了从\textbf{代码到数据的身份扭曲}。这是对图灵机架构基本原则的违反，尽管在通用图灵机的意义下是合法的。
	
	\textbf{拷贝不可能性作为根源。} 这一扭曲的深层原因是第4节所揭示的拷贝不可能性。经典模型无法执行从虚部到实部的拷贝——它无法将控制信息转换为数据，因为它的操作语义根本不包含``虚部''和``控制信息''这些概念。因此，经典NP定义别无选择，只能将$c$直接写在实部纸带上，从而必然造成身份扭曲。这不是一个可以修补的技术疏忽，而是经典模型操作语义边界的必然表现。
	
	\textbf{不可公度性的丢失。} 这一身份扭曲不仅违反了代码与数据分离的原则，更导致了一个同等严重的后果：\textbf{见证$c$的各个比特之间原本的不可公度性被完全抹除。}
	
	在虚部中，$c$的第$i$个比特对应于在第$i$个非确定性选择点上是选择生成元$+\sqrt{p_i}$还是$-\sqrt{p_i}$。这些生成元（如$\sqrt{p_2}$和$\sqrt{p_3}$）在基域$\mathbb{Q}$上不可公度——它们之间不存在任何非零有理线性关系。这意味着，选择了$\sqrt{p_2}$所标记的路径（第一个分支的激活路径）与选择了$\sqrt{p_3}$所标记的路径（第二个分支的激活路径），在代数上是不可相互归约的。这正是非确定性选择独立性的数学基础——不同分支步骤之间的独立性由对应生成元的不可公度性严格保证。
	
	然而，当这些选择被编码为数据比特$0$和$1$后，比特之间不再保留任何代数独立性关系。$0$和$1$只是两个不同的符号，它们在基域内完全可以相互转化——交换$0$和$1$的语义标签，机器的操作轨迹不变（这正是语法不变性引理的结论，见第6.1节）。因此，经典模型无法从数学上区分``两个真正独立的非确定性选择''与``两个仅在语法标签上不同的确定性选项''。\textbf{非确定性选择的独立性在概念层面被彻底抹除。}
	
	这解释了为什么经典框架从未正面定义``非确定性选择的独立性''——不是疏忽，而是因为其定义本身就将保证独立性的代数结构（不可公度性）丢失了。
	
	\subsection{第二重错误：对偶性被存在量词消除}
	
	第一重错误已经足够严重，但更深刻的问题发生在下一步。
	
	在非确定性分支中，每一次选择都是在两个\textbf{对偶的生成元}之间做出的。在CBTM框架中，这两个生成元是$\alpha$和$\beta$，它们在基域$\mathbb{F}_2$上不可公度，且由伽罗瓦自同构$\sigma: \alpha \leftrightarrow \beta$精确地相互映射。在IVM框架中，生成元为$\sqrt{p_i}$和$-\sqrt{p_i}$，它们满足相同的代数方程$x^2 - p_i = 0$，且由符号翻转构成对偶。一个选项之所以有意义，恰恰是因为它与另一个选项的对立——非确定性选择的本质结构就是这种\textbf{对偶性}。
	
	经典NP定义通过以下操作消除了这一对偶性：
	\begin{enumerate}
		\item \textbf{编码为比特串}：原本的``选择$\alpha$还是选择$\beta$''被扁平化为``位$i$是0还是1''。0和1在语法上没有对偶关系——它们只是两个不同的符号，不存在伽罗瓦自同构将它们精确地相互映射。
		\item \textbf{存在量化}：量词``$\exists c$''只断言存在某个比特串使验证器接受。它不记录这个比特串是在哪些对偶选项中做选择的结果，不保留那个``未被选择的另一条路径''的任何痕迹。对偶性——这个非确定性结构中最核心的代数属性——在量词的包裹下被彻底删除。
	\end{enumerate}
	
	\textbf{编码的结构性代价。} 需要特别指出的是，对偶性的丢失不是偶然的，而是从虚部到实部的编码过程的\textbf{结构性必然}。在虚部中，生成元$+\sqrt{p_i}$与$-\sqrt{p_i}$满足相同的代数方程$x^2-p_i=0$，由伽罗瓦自同构$\sigma_i: \sqrt{p_i} \leftrightarrow -\sqrt{p_i}$精确地相互映射。这两个选项之间的对偶关系是一个代数事实。然而，比特$0$和$1$之间不存在任何类似的代数映射——没有任何自同构将$0$和$1$精确地相互交换并保持基域不动。因此，任何将虚部中的对偶选择编码为比特串的操作，都必然丢失对偶性。这不是编码方案的选择问题，而是两种语言表达能力的原则性差异。
	
	\textbf{拷贝不可能性作为根源：结构性的必然丢失。} 从虚部到实部的编码过程中，对偶性的丢失是\textbf{结构性的必然}。经典模型的语言只能表达离散的比特，而虚部中的对偶结构是连续的代数关系。任何将后者编码为前者的操作，都必然丢失对偶性。因此，经典NP定义中对偶性的缺失，根源在于它试图用经典模型的语言（只能表达比特）去描述一个本质上超出该语言表达能力（需要虚部空间和伽罗瓦群）的结构。
	
	结果是：\textbf{经典框架所说的``非确定性''，是一个去掉了非确定性本质（对偶性）之后的残缺概念。} 经典NP定义不仅没有正面定义对偶性，而且在从虚部到实部的编码过程中，必然地、结构性地消除了它。
	
	\subsection{第三重错误：语义外包导致概念残缺}
	
	综合前两重错误，经典NP定义的根本问题可以概括为：
	
	\begin{quote}
		\textbf{它用语法层面的``存在多个五元组''来定义非确定性，但非确定性的本质——分支选项之间的不可公度性及其对偶结构——是代数语义层面的属性。经典确定性图灵机的内部语义无法表达不可公度性（这是不可公度性元定理的结论），且无法执行从虚部到实部的拷贝（这是投影约束公理的推论）。更为致命的是，经典NP定义在将见证$c$写在实部纸带上的过程中，\textbf{已经丧失了不可公度性}（第一重错误的后果）和\textbf{对偶性}（第二重错误的后果）。当独立性和结构都已丢失后，经典模型除了将非确定性的语义本质外包给外部量词和外部谓词，别无选择。\textbf{第三重错误是前两重错误在语义层面的必然后果。}}
	\end{quote}
	
	这一``语法定义、语义外包、拷贝不可能、概念扭曲''的模式造成了三个严重后果：
	\begin{enumerate}
		\item \textbf{非确定性选择的独立性从未被正面定义}：在经典框架内，你无法回答``为什么这两个后继是真正的非确定性选择''这一问题，因为回答它需要不可公度性和对偶性的概念，而经典操作语义无法内在地表达前者，经典NP定义又在编码过程中必然地丢失了两者。
		\item \textbf{三大障碍正是利用了这一语义空洞和概念扭曲}：相对化障碍的根源在于谕示可以替机器完成“猜测虚部控制信息并拷贝到实部”这一不可能操作，自然证明障碍基于无法触及不可公度性的组合性质，代数化障碍只能回答多项式等式问题而无法感知符号正负号——它们之所以能长期存在，正是因为经典模型缺乏正面定义独立性所需的语义资源，也无法区分激活与不激活两种对偶状态。
		\item \textbf{P vs NP问题在经典框架内成为``语言性的''困难}：问题不在于计算本身，而在于用以描述计算的语言缺乏必要的概念词汇，且经典NP定义在操作语义上是不自足的。经典框架中的``非确定性''，是一个在定义层面就被双重残缺化的概念——既在NTM层面缺失不可公度性的正面定义，又在NP层面因编码必然丢失不可公度性和对偶性。
	\end{enumerate}
	
	\subsection{修正的方向}
	
	上述诊断直接指向了修正的方向。要使得P vs NP问题获得合法的表述和可证的解答，我们需要一个计算模型，它满足以下要求：
	\begin{enumerate}
		\item \textbf{不可公度性内建}：将不可公度性作为操作语义的内在属性，而非在函数指针被扭曲为数据的过程中将其丢失。
		\item \textbf{消除拷贝操作}：让非确定性选择的函数指针地址常驻于控制层（虚部），从根本上消除``从虚部到实部拷贝''这一不可能操作的必要性。
		\item \textbf{对偶性常驻}：保留非确定性选择的对偶结构，而非将其编码为无对偶关系的比特串并用量词消除。
		\item \textbf{语义自足}：非确定性由机器内部的代数机制来定义和执行，而非由外部量词``$\exists c$''来断言。
	\end{enumerate}
	
	CBTM/IVM框架正是按照这些要求构建的\cite{cbtmarxiv,rbtmarxiv}。这解释了为什么新框架能够完成经典框架无法完成的证明：它不是简单地``重新解释''非确定性，而是\textbf{补全了经典定义中长期缺失的语义前提}，\textbf{消除了经典定义中不可能操作的依赖}，并\textbf{修正了经典定义中由此必然导致的概念扭曲}。
	
	\subsection{语义缺失的必然技术后果：三大障碍的根源} \label{sec:barriers-root}
	
	前述三重概念错误的诊断并非孤立的哲学批评，它们直接解释了为什么相对化、自然证明和代数化三大障碍能够在经典框架中长期存在。事实上，这三大障碍正是语义缺失、操作不自足和概念扭曲在不同技术层面的必然表现。
	
	\textbf{相对化障碍与操作不自足。} 在经典NP的定义中，见证$c$是由外部量词“$\exists c$”所断言的——外部观察者通过猜测的方式，获得了原本位于虚部的控制信息（即非确定性选择的轨迹），然后将其拷贝到实部纸带上，使之成为验证器$V$可读取的数据$c$。这正是本文第4节所诊断的那个不可实现的拷贝操作。从CBTM架构来看，虚部与实部分属控制层与数据层，二者之间存在严格的投影约束：经典机器无法感知虚部，更无法将虚部信息跨层次拷贝到实部。然而，在相对化世界中，谕示可以被暗中赋予填补这一语义空洞的超能力——替机器完成“猜测虚部控制信息并拷贝到实部”这一经典模型自身不可能执行的操作。当谕示免费承担了这一角色，$\mathbf{P}^A = \mathbf{NP}^A$ 便出现了。因此，相对化障碍的深层根源在于：谕示被暗中赋予了经典机器所不具备的外部资源，这正是经典NP定义操作不自足性的必然表现。这表明相对化世界中$\mathbf{P}^A=\mathbf{NP}^A$的构造依赖于填补操作不自足的外部资源，并不构成对真实计算世界中P与NP关系的反证。
	
	\textbf{自然证明障碍与代数概念的超组合性。} 自然证明要求下界证明所基于的组合性质对大多数布尔函数成立（“大尺度”条件）。然而，非确定性选择的独立性——即不同分支选项在代数上不可相互归约——本质上是一种代数概念（不可公度性），而非组合概念。经典NP定义在将见证$c$编码为比特串时，已经丢失了这种代数独立性（第一重错误）；经典框架因而只能用组合语言去表述一个其本质是代数的问题。自然证明之所以成为障碍，正是因为它只能触及布尔函数的组合属性，而无法触及“真正独立的非确定性选择”所需的代数结构。我们将在后续论文中严格论证，任何能够分离P与NP的不变量，其定义域必然严格限于NP语言，从而结构性规避自然证明的“大尺度”条件。
	
	\textbf{代数化障碍与对偶性的丧失。} 代数化框架将谕示替换为低次多项式扩展。然而，代数化谕示所能回答的所有问题，归根结底都是关于多项式等式的问题——它无法感知符号的正负号。在非确定性分支中，两个对偶选项分别对应于生成元的正负号（如 $+\sqrt{p_i}$ 与 $-\sqrt{p_i}$），这正是被经典NP定义用存在量词消除掉的对偶性（第二重错误）。代数化谕示对这种符号信息天然“失明”。因此，如果P与NP的差异恰好寓于这种对偶性之中——即非确定性分支的触发依赖于对偶生成元的选择，而确定性计算无法在不借助该信息的情况下等价模拟——那么任何代数化框架都无法触及这一差异。可以严格证明，基于伽罗瓦自同构敏感性的分离论证是谕示无关的，因而完全不受代数化障碍的约束。
	
	进一步看，代数化障碍实质上是相对化障碍在多项式域上的延续。正因如此，在 CBTM/IVM 框架中为抵抗相对化而设的投影约束机制——将一切外部相应锁定在实部——同时也就剥夺了代数化谕示触及对偶性的任何可能。破除两者的，是同一种语义隔离。
	
	综上，三大障碍并非计算本身的绝对限制，而是经典NP定义在语义缺失、操作不自足和概念扭曲的框架下，任何纯语法论证都难以避免的外部性困境。它们的存在恰恰印证了本文的核心诊断：经典语言缺乏正面定义非确定性独立性和结构所需的概念资源（不可公度性与对偶性），且其定义本身依赖于一个不可实现的操作。要克服这些障碍，必须将不可公度性内建为操作语义的内在属性，让对偶性常驻虚部，并从根本上消除拷贝操作的必要性——这正是后续论文中CBTM/IVM框架的任务。
	
	\section{不可公度性元定理}
	
	第3-5节的诊断指向一个关键问题：经典图灵机是否真的无法内在地表达不可公度性？本节给出这一命题的严格数学证明。
	
	\subsection{语法不变性：经典操作语义的核心}
	
	经典图灵机的转移函数$\delta$仅依赖符号的标识（形状），而非符号可能携带的任何外部语义。这一事实可以形式化为以下引理。
	
	\begin{definition}[经典图灵机]
		一台经典（确定性）图灵机是一个七元组$M = (Q, \Sigma, \Gamma, \delta, q_0, q_{\text{accept}}, q_{\text{reject}})$，其中$Q$是有限状态集，$\Gamma$是有限纸带字母表，$\delta: Q \times \Gamma \to Q \times \Gamma \times \{L, R\}$是转移函数。字母表$\Gamma$包含一个特殊的空白符$B$，初始格局中除输入字符串外纸带上所有单元均为空白符。
	\end{definition}
	
	转移函数$\delta$的核心特征是：其输入中的符号参数仅仅是$\Gamma$中的一个\textbf{标识}。$\delta$无法区分两个具有不同外部语义但形式上可互换的符号，它只能根据当前状态和当前符号的标识来决定下一步动作。
	
	\begin{lemma}[语法不变性引理]\label{lem:syntactic-invariance}
		设$M = (Q, \Gamma, \delta, q_0, q_{\text{accept}}, q_{\text{reject}})$为一台经典图灵机，其中$\Gamma$包含空白符$B$。设$\pi: \Gamma \to \Gamma$为字母表上的置换（双射），满足$\pi(B) = B$（即置换保持空白符不变）。构造新机器$M_\pi = (Q, \Gamma, \delta_\pi, q_0, q_{\text{accept}}, q_{\text{reject}})$，其中对于所有$q \in Q, c \in \Gamma$：
		\[
		\delta_\pi(q, \pi(c)) = (q', \pi(c'), D) \iff \delta(q, c) = (q', c', D).
		\]
		则$M$与$M_\pi$结构同构：若$M$在输入$w = w_1 w_2 \ldots w_n$上经历格局序列$C_0, C_1, C_2, \ldots$，则$M_\pi$在输入$\pi(w) = \pi(w_1) \pi(w_2) \ldots \pi(w_n)$上经历格局序列$\pi(C_0), \pi(C_1), \pi(C_2), \ldots$，其中$\pi(C)$表示将格局$C$纸带上所有符号$c$替换为$\pi(c)$所得到的格局。
	\end{lemma}
	\begin{proof}
		对计算步数进行归纳。条件$\pi(B)=B$保证$M_\pi$的初始格局是合法的（非输入部分仍为空白符）。归纳基础：初始格局$C_0$与$\pi(C_0)$对应。归纳步骤：若第$t$步$M$在格局$C_t$处读取符号$c$并转移到$C_{t+1}$，则$M_\pi$在格局$\pi(C_t)$处读取符号$\pi(c)$，根据$\delta_\pi$的定义，执行相同的动作并转移到$\pi(C_{t+1})$。
	\end{proof}
	
	语法不变性引理揭示了一个基本事实：只要不置换空白符，图灵机的行为在数据符号重标记下是完全对称的。将所有的``0''换成``1''、所有的``1''换成``0''，机器的操作轨迹不变，仅仅是符号的标签变了。
	
	\subsection{内在属性与置换不变性}
	
	\begin{definition}[内在属性]
		设$\mathfrak{P}$是机器$M$在输入$w$上的一个属性。称$\mathfrak{P}$是图灵机操作语义的\textbf{内在属性}，如果它完全由机器的转移规则和符号标识决定，而不依赖于外部观察者对符号的语义赋值。形式化地，$\mathfrak{P}$是内在属性当且仅当对于字母表上任意保持空白符不变的置换$\pi: \Gamma \to \Gamma$，$M$在输入$w$上具有$\mathfrak{P}$当且仅当$M_\pi$在输入$\pi(w)$上具有$\mathfrak{P}$。
	\end{definition}
	
	这一形式化刻画基于经典图灵机操作语义的核心特征：转移函数$\delta$仅依赖当前状态和读写头下符号的标识，而不依赖任何外部语义。因此，机器的全部行为完全由状态集$Q$、字母表$\Gamma$、以及基于这些集合定义的转移规则$\delta$所决定。在此设定下，任何能够被机器``内在地''区分或响应的属性，必须在字母表$\Gamma$的重命名（即置换）下保持不变，因为重命名仅仅改变了符号的标签，而没有改变状态结构和转移规则。内在属性的置换不变性因此穷尽了基于语法规则本身所能区分的所有属性，是经典图灵机语法本质的直接体现。
	
	例如，``机器在输入$w$上停机''是一个内在属性——无论我们将符号解释成什么数学对象，停机的判断只取决于操作规则本身。以下推论是语法不变性引理的直接结果。
	
	\begin{corollary}[内在属性的置换不变性]\label{cor:intrinsic}
		若属性$\mathfrak{P}$是机器$M$操作语义的内在属性，则对于任意保持空白符不变的置换$\pi: \Gamma \to \Gamma$，$M$在输入$w$上具有$\mathfrak{P}$当且仅当$M_\pi$在输入$\pi(w)$上具有$\mathfrak{P}$。
	\end{corollary}
	
	其逆否命题是本文证明的核心逻辑：\textbf{如果一个属性在某个置换下发生了改变，那么该属性必然依赖于外部语义赋值，不可能是操作语义的内在属性。}
	
	\subsection{编码不可公度性}
	
	为了严格地讨论``机器是否能处理不可公度元''，我们必须显式地引入外部语义。这通过一个编码函数来完成。
	
	\begin{definition}[代数不可公度性]
		两个代数数$\gamma, \delta$在有理数域$\mathbb{Q}$上\textbf{代数不可公度}，当且仅当它们在$\mathbb{Q}$上线性无关，即不存在不全为零的有理数$c_1, c_2$使得$c_1\gamma + c_2\delta = 0$。这等价于要求$\gamma$与$\delta$皆非零，且其比值$\gamma/\delta$为无理数。
	\end{definition}
	
	需要指出，对于两个元素，``不可公度''与``线性无关''是等价的。本文定理直接处理两个元素的情形，因此使用``不可公度性''一词。对于多个元素，``线性无关性''是``不可公度性''的自然推广。
	
	这一定义确保了零元素不参与不可公度性的判定：包含零的集合天然是线性相关的（可公度）。
	
	\begin{definition}[代数编码与编码不可公度性]
		设$\varepsilon_0: \Gamma \setminus \{B\} \to \mathbb{A}$是一个为每个数据符号指定一个代数数的函数。将其扩展为$\varepsilon: \Gamma^* \to \mathbb{A}$，使得对于输入$w \in \Gamma^*$，$\varepsilon(w) = \{\varepsilon_0(c) : c \text{ 是 } w \text{ 中出现的符号}\}$。称机器$M$在输入$w \in \Gamma^*$上\textbf{关联$\varepsilon$不可公度元}，如果$\varepsilon(w)$中包含两个非零代数数$\gamma, \delta$，它们在有理数域$\mathbb{Q}$上代数不可公度。
	\end{definition}
	
	注意，此处$\varepsilon(w)$被定义为符号值构成的集合而非多重集，因为线性相关性仅取决于元素的存在性，与重复次数无关。更重要的是，这一定义将``关联不可公度元''界定为一个\textbf{仅依赖于输入字符串$w$和编码$\varepsilon$、而与机器$M$的计算行为无关}的外部属性。这意味着我们考察的是一个比``机器在计算中利用了不可公度元''更弱的条件，即仅仅是``输入中包含了不可公度元''。如果经典模型连这种最弱意义上的不可公度性都无法内在地表达，那么任何更强的、涉及机器行为的定义就更加不可能。这一设计使定理的结论更强而非更弱。
	
	\subsection{主要定理及其完整证明}
	
	为聚焦于操作语义的本质而非编码细节，下文考虑具备足够表达能力的字母表。在P vs.\ NP等标准复杂性分析语境下，通常默认允许多带图灵机或多符号字母表，因为它们不影响多项式时间等价类的定义。因此，本文的结论适用于这一标准分析语境下的图灵机变体。
	
	\begin{theorem}[不可公度性不能被内在化]\label{thm:main}
		设经典图灵机的字母表$\Gamma$包含空白符$B$及至少四个不同于$B$的数据符号。则不存在一个操作语义的内在属性$\mathfrak{P}$，使得对于所有满足以下条件的编码函数$\varepsilon$——即$\varepsilon$将每个数据符号独立地映射到某个代数数，且将任意两个指定的数据符号映射到非零有理数$a$和$b$（其中$a/b \in \mathbb{Q} \setminus \{0\}$）——和所有输入$w \in \Gamma^*$（其中$w$不包含空白符$B$），$\mathfrak{P}(M, w)$成立当且仅当$M$在$w$上关联$\varepsilon$不可公度元。等价地，经典图灵机无法通过其操作语义唯一地``锚定''不可公度性。
	\end{theorem}
	
	\noindent\textbf{注：}即使限定于三符号$\{0,1,B\}$模型，通过标准配对函数可将$k$个符号压缩为$\lceil\log_2 k\rceil$比特，所得多带模拟在多项式时间内等价，故定理结论适用于$\mathbf{P}/\mathbf{NP}$的标准分析语境。
	
	\begin{proof}
		设$\Gamma$的数据符号包含$\{0, 1, \alpha, \beta\}$，且它们均不同于空白符$B$。作为构造的一部分，我们选取数据符号$0$和$1$作为定理条件中所指的``两个指定符号''。我们采用反证法。
		
		\textbf{假设}：存在一个内在属性$\mathfrak{P}$，它能够精确地捕获``关联不可公度元''这一概念。即，对于任何满足条件的编码$\varepsilon$和任何由数据符号组成的输入$w$，$\mathfrak{P}(M, w)$为真当且仅当$M$在$w$上关联$\varepsilon$不可公度元。
		
		\textbf{步骤1：固定一个满足条件的编码$\varepsilon$。}
		定义编码$\varepsilon_0: \Gamma \setminus \{B\} \to \mathbb{A}$，它将数据符号映射为：
		\[
		0 \mapsto 1,\quad 1 \mapsto 2,\quad \alpha \mapsto \sqrt{2},\quad \beta \mapsto \sqrt{3},
		\]
		并扩展为$\varepsilon$。\footnote{这里有意将符号``0''和``1''映射到与自身标签不同的非零有理数（1和2），以强调编码函数仅依赖符号的标识（标签）而非符号名称的数值暗示。经典图灵机的操作语义只认识符号的标签，而不认识其被赋予的数值，这正是定理所要揭示的局限。}此编码满足定理条件：指定的两个符号被映射到非零有理数，且比值为有理数。
		
		\textbf{步骤2：构造具体输入并应用假设。}
		考虑任意一台图灵机$M$在输入$w = \alpha\beta$上的行为。根据$\varepsilon$，$\varepsilon(\alpha\beta) = \{\sqrt{2}, \sqrt{3}\}$，两者皆非零且在$\mathbb{Q}$上线性无关（不可公度）。因此，$M$在$w$上关联$\varepsilon$不可公度元。由假设，内在属性$\mathfrak{P}(M, \alpha\beta)$必须为真。
		
		\textbf{步骤3：构造符号置换$\pi$和置换后的机器$M_\pi$。}
		定义置换$\pi: \Gamma \to \Gamma$，保持空白符不变（$\pi(B)=B$），并在数据符号上的作用为：
		\[
		\pi(0) = \alpha,\quad \pi(1) = \beta,\quad \pi(\alpha) = 0,\quad \pi(\beta) = 1.
		\]
		这是一个保持空白符不变的双射。根据引理\ref{lem:syntactic-invariance}，构造$M_\pi$，它与$M$结构同构。
		
		\textbf{步骤4：应用内在属性的置换不变性。}
		由于$\mathfrak{P}$被假设为内在属性，根据推论\ref{cor:intrinsic}，它在符号置换下必须保持不变。因此，既然$\mathfrak{P}(M, \alpha\beta)$为真，那么$\mathfrak{P}(M_\pi, \pi(\alpha\beta)) = \mathfrak{P}(M_\pi, 01)$也必须为真。
		
		\textbf{步骤5：在同一编码$\varepsilon$下检验矛盾。}
		对于$M_\pi$的输入$01$，我们仍使用步骤1中固定的编码$\varepsilon$。根据$\varepsilon$，$\varepsilon(0 1) = \{1, 2\}$，两者皆为非零有理数，它们在$\mathbb{Q}$上线性相关（可公度），且不包含任何非零的不可公度元对。因此，$M_\pi$在输入$01$上并不关联$\varepsilon$不可公度元。
		
		\textbf{步骤6：导出矛盾。}
		根据步骤4，$\mathfrak{P}(M_\pi, 01)$为真。但根据步骤5和反证假设，$\mathfrak{P}(M_\pi, 01)$应为假（因为它不关联不可公度元）。这是一个矛盾。
		
		这个矛盾表明，我们最初的假设——存在一个内在属性$\mathfrak{P}$能够独立于编码地捕获关联不可公度元——是错误的。需要特别指出的是，本定理中的``关联不可公度元''被定义为纯粹的外部属性（仅取决于输入和编码），这一设计恰恰强化了结论：如果经典模型连这种最弱意义上的不可公度性都无法内在地锚定，那么任何更强的、涉及机器行为的定义就更加不可能。我们的证明用形式化手段锚定了这一直觉：任何内在属性都无法在符号置换下保持与这种外部属性的等价对应。
	\end{proof}
	
	该证明的美感在于其极简性。它不依赖于任何关于图灵机计算能力的假设，不涉及多项式时间、非确定性或任何复杂性类的概念。它仅使用了图灵机最基本的语法性质——转移函数仅依赖符号标识——以及一个精巧的符号置换构造。
	
	\subsection{关于字母表大小的边界情况}
	
	定理\ref{thm:main}要求字母表$\Gamma$包含空白符$B$及至少四个不同于$B$的数据符号，即$|\Gamma| \ge 5$。对于标准的、仅含$\{0,1,B\}$的三符号图灵机模型，可通过标准的配对编码方案将$k$个符号压缩为$\lceil \log_2 k \rceil$个比特，从而在多项式时间等价的扩展模型中应用本定理。即使在通过配对编码扩展的模型中，定理的置换不变性论证框架依然适用。核心要点在于：问题的根源在于转移函数的语法性，而非字母表的具体大小。因此，本文的结论适用于复杂性理论的标准分析语境。
	
	\subsection{元定理的意义}
	
	元定理的直接推论是：\textbf{经典图灵机的操作语义本身永远无法``感知''不可公度性。}这不是在``算法可判定性''的弱意义上（经典图灵机可以在适当编码下判定不可公度性），而是在``操作语义内蕴性''的强意义上（它不能内在地锚定不可公度性）。无论编码多么精巧，经典图灵机的操作语义本身永远无法区分``这个符号编码了$\sqrt{2}$''和``这个符号编码了$1$''——因为符号可以被置换，而机器行为在置换下保持不变。
	
	因此，外部谓词$IsWitness_L(c)$在经典框架内是一个\textbf{语义空洞}。它指向了一个经典语言无法表达的核心概念。NP定义将非确定性的本质外包给外部量词和外部谓词，并非疏忽，而是其语言局限的必然表现——经典模型缺乏表达不可公度性所需的概念资源。
	
	\subsection{代数封闭性的视角}
	
	语法不变性（置换论证）从语法层面证明了不可公度性不是内在属性。代数封闭性从另一个互补的角度强化了这一结论。
	
	经典图灵机的操作语义，其代数表达力被限制在编码所选择的基域之内。以有理数域$\mathbb{Q}$为例：它对加法、减法、乘法、除法（非零）封闭。从有理数出发，通过有限步有理运算，永远无法``跳''到无理代数数。因此，如果编码函数将符号映射到有理数，机器通过转移函数能够操作和生成的任何信息，都永远落在$\mathbb{Q}$内。
	
	而不可公度性（如$\sqrt{2}$和$\sqrt{3}$的关系）是一种跨越基域边界的属性，它只有在扩张域$\mathbb{Q}(\sqrt{2}, \sqrt{3})$中才有定义。经典图灵机的操作语义在原则上就无法触及它。置换论证与代数封闭性论证从不同角度指向同一结论：不可公度性不可能是操作语义的内在属性。
	
	\section{``聪明算法''假设中的语义错位}
	
	\subsection{``聪明算法''假设实际要求什么}
	
	当我们假设存在一个``聪明算法''——即一个$\mathbf{P}$算法，能够在多项式时间内判定某个$\mathbf{NP}$完全问题（如子集和）——时，我们实际上是在假设什么？
	
	从经典模型的角度看，$\mathbf{P}$算法是经典模型内部的一个对象。它的操作语义完全由转移函数定义——每一步可以抽象为状态空间中的仿射变换$X \leftarrow X + C$。$\mathbf{P}$算法只能``理解''和``表达''这种内部的、操作性的语义。它无法产生或利用不可公度生成元，因为它的操作语义在原则上无法触及跨越基域边界的代数结构。更为根本的是，\textbf{它无法执行从虚部到实部的拷贝操作}——经典P算法只能操作实部数据，它的操作语义中根本没有``虚部''的概念。
	
	然而，$\mathbf{P}$算法被要求判定的$\mathbf{NP}$完全问题，其非确定性验证过程的本质恰恰需要不可公度性和对偶性来保证。见证$c$的每一个比特代表一次独立的函数指针选择——在CBTM/IVM框架中，这是在对偶生成元$\alpha$与$\beta$（或$\sqrt{p_i}$与$-\sqrt{p_i}$）之间做出的选择。而且，\textbf{NP验证的定义本身依赖于从虚部到实部的拷贝结果}——经典定义中的见证$c$正是这一拷贝的产物。经典P算法无法执行这一拷贝，却要等价地实现依赖于这一拷贝结果的验证过程。
	
	这构成了一个语义层面的困境：\textbf{聪明算法假设要求一个受限于经典模型内部语义的对象（无法表达不可公度性，无法执行虚部到实部拷贝，且丧失了对偶性），去等价地实现一个依赖于不可公度性、依赖于该拷贝结果、且依赖于对偶性的外部语义。}这正是在要求用不存在不可公度性、无法执行必要操作、且丧失了对偶性的模型，去实现不可公度性、执行不可能操作、并复现对偶性——语义层面和操作层面的双重资源错位。
	
	\subsection{语义困境：实现不可表达之物}
	
	基于以上分析，我们可以精确地刻画聪明算法假设中的语义错位。
	
	它不是可从$\mathbf{P}=\mathbf{NP}$推导出形式逻辑矛盾（如$0=1$）的意义上的``悖论''。它是一种更微妙但同样根本性的困境：\textbf{系统被要求去完成一项其在原则上无法完成的任务。}
	
	具体而言，经典图灵机本质上在有理数域上进行仿射计算——每一步操作可以抽象为$X \leftarrow X + C$。它无法触及无理代数数的不可公度性，无法保留非确定性选择的对偶结构，更无法执行从虚部到实部的拷贝操作。而非确定性选择的独立性恰恰需要不可公度性来保证，其结构恰恰是对偶的，其定义本身恰恰依赖于那个不可执行的拷贝操作。这不是计算能力的不足，而是概念资源和操作能力的双重缺失。更有趣的是，如果$\mathbf{P}=\mathbf{NP}$，那么NP定义中的外部量词``$\exists c$''和外部谓词$IsWitness(c)$就被证明是冗余的，NP概念自我消解，命题本身失去语义基础。
	
	\section{哥德尔与塔斯基的启示}
	
	\subsection{哥德尔不完备定理}
	
	1931年，哥德尔证明了两个震撼数学基础的定理\cite{Godel1931}。第一不完备定理：任何包含算术的、递归可公理化的、一致的形式系统，都存在一个真命题，它在系统内部不可证。第二不完备定理：这样的系统无法证明自身的一致性。哥德尔定理的核心洞见是：一个足够强大的形式系统的语义（真理性）不能由系统内部的语法操作完全定义。系统无法``跳出自己''去确证自身的无矛盾性。语义必须从外部引入。
	
	\subsection{塔斯基的真不可定义定理}
	
	塔斯基\cite{Tarski1935}进一步阐明：在足够丰富的语言中，真概念不能在语言内部定义。如果一个人试图在对象语言内部定义一个谓词$T(x)$，使得$T(\ulcorner\phi\urcorner)$成立当且仅当$\phi$为真，那么他将不可避免地遭遇说谎者悖论。真概念必须位于比对象语言更高的元语言层次中。这意味着\textbf{语义总是超越语法}。一个形式系统的语义不能在该系统内部被完全捕获；它需要外部概念资源的注入。
	
	\subsection{与经典图灵机的类比}
	
	将哥德尔和塔斯基的洞见应用于计算理论，我们得到以下类比。
	
	经典图灵机是一个纯语法的计算模型。正如哥德尔证明了算术系统无法在内部证明自身的一致性，不可公度性元定理证明了经典图灵机无法在内部（在操作语义的强意义上）表达不可公度性。正如塔斯基证明了真概念不能在对象语言内部定义，我们揭示了非确定性选择的独立性不能在经典操作语义内部定义。
	
	在这两种情况下，困境的根源都不是技术性的，而是\textbf{语言性的}。哥德尔的不可证命题在系统内部为真，但其证明需要比原系统更强的元理论工具（如Gentzen使用超限归纳法证明了算术的一致性）。同样，$\mathbf{P} \neq \mathbf{NP}$在经典框架内为真（作为集合论命题），但其证明需要比经典模型更丰富的语义框架——一个能够内在地表达不可公度性、保留对偶性、并从根本上消除拷贝操作依赖的计算模型。
	
	\subsection{类比的边界}
	
	然而，这一类比有其边界，需要精确界定。
	
	哥德尔的不完备定理是关于\textbf{可证性}的：一致性不可证。我们的元定理是关于\textbf{可表达性}的：不可公度性不可内在表达。两者涉及不同的逻辑范畴。哥德尔的证明使用了自指性构造（哥德尔编码），而我们的证明使用了符号置换和语义重赋值。
	
	更重要的是，哥德尔不完备定理是经典数学框架内部的严格定理。同样，我们的不可公度性元定理也是在经典数学框架内部被严格证明的。由此出发得出的哲学诊断，即经典框架缺乏证明所需的概念资源且经典NP定义在操作上不自足，并不等同于断言该问题没有真值或无法在更丰富的框架中被解决。本文在这一问题上的立场是审慎的：我们主张的是，经典框架的语言不足以正面定义非确定性选择的独立性且无法执行必要的操作，而非该问题在绝对意义上是``不合法的''。
	
	\section{与BSS模型的比较}
	
	\subsection{BSS模型的策略：扩展计算域}
	
	1989年，Blum、Shub和Smale\cite{Blum1989}提出了BSS模型，将计算复杂性理论从离散的图灵机世界推广到连续的实数世界。BSS模型的核心特征是允许每个寄存器存储任意实数，并在单步内执行实数上的精确算术运算（加、减、乘、除）和比较。
	
	这一推广具有深远的理论动机。许多连续数学中的算法（如牛顿法、线性规划）涉及实数运算，经典图灵机模型无法精确刻画其计算复杂性。BSS模型试图通过扩展计算域——从有限字母表到整个实数域——来捕捉连续计算的本质。
	
	然而，这种扩展偏离了丘奇–图灵论题的离散预设：BSS 模型将``实数''视为可一步读取和运算的原子对象，相当于在每一步隐式注入了无穷位信息。这使得某些在经典图灵机下不可判的问题（如判定一个实数是否为整数）在 BSS 模型下变为可判定，但这一``增益''来源于对实数原子的理想化假定，而非模型本身超越了图灵可计算性的界限。
	
	\subsection{为什么BSS模型未能解决经典P vs.\ NP}
	
	BSS模型中的$\mathbf{P}_{\text{BSS}}$与$\mathbf{NP}_{\text{BSS}}$的关系也是开放的。尽管BSS模型提供了丰富的代数结构，它并未解决经典的P vs.\ NP问题。原因可以从三个层面来理解。
	
	\textbf{第一，计算域的扩展导致复杂性类的不对应。} BSS模型建立在实数域上，而经典P vs.\ NP问题定义在有限字母表上。BSS模型通过假设实数可以被精确存储和操作，跳过了``从离散到连续''的编码鸿沟。这种跳跃偏离了丘奇–图灵论题的离散预设，使得BSS的复杂性类与经典复杂性类之间缺乏直接的对应。
	
	\textbf{第二，拷贝操作的不可能性未被解决。} BSS模型虽然允许实数运算，但仍未解决经典NP定义中隐含的``虚部到实部拷贝''问题。BSS模型中的``猜测''仍然由外部观察者提供——机器本身不知道如何从无限多种可能性中选择正确的实数见证。经典NP定义的操作不自足性在BSS模型中依然存在。
	
	\textbf{第三，不可公度性的来源不同。} 在BSS模型中，实数之间当然存在不可公度性（如$1$与$\pi$在$\mathbb{Q}$上线性无关）。但这种不可公度性是\textbf{外部赋予}的——它来自实数域本身的代数结构，而非由机器的操作语义所产生。BSS模型将不可公度性作为计算域的一个既定属性，而非操作语义的内在组成部分。更关键的是，BSS模型并未将不可公度性与非确定性分支的独立性建立关联——它没有将生成元的对偶性与非确定性选择的二元性对应起来。
	
	\subsection{另一条路径：内建语义而非扩展域}
	
	CBTM/IVM框架采取了与BSS模型互补但本质不同的路径。
	
	\textbf{BSS模型}的策略是：\textbf{扩展计算域}——从有限字母表扩展到实数域，让机器操作更丰富的数学对象。代价是偏离了丘奇–图灵论题的离散预设，与经典P vs.\ NP问题的语境不再直接对应，且未解决拷贝操作的不可能性问题。
	
	\textbf{CBTM/IVM框架}的策略是：\textbf{内建语义}——在保持离散计算的前提下，将不可公度性和对偶性作为操作语义的内在属性，并从根本上消除拷贝操作的必要性。CBTM的纸带符号取自四元有限域$\mathbb{F}_4$，与经典图灵机在计算能力上完全等价（参数化等价定理）。但通过虚部机制和分支触发公理，CBTM将不可公度性从外部语义赋值转化为操作语义的内部结构。当虚部$b=1$时，机器\textbf{强制}分裂为两条路径，分别绑定于不可公度生成元$\alpha$和$\beta$——这不是因为符号``恰好''具有不可公度的外部解释，而是因为操作语义本身要求如此。同时，函数指针地址常驻虚部，不再需要拷贝到实部；$\alpha$与$\beta$的对偶关系由伽罗瓦自同构精确刻画，对偶性成为操作语义的显式结构。
	
	\subsection{启示}
	
	BSS模型与CBTM/IVM框架的比较揭示了一个关键的方法论洞见：\textbf{单纯地将符号提升为代数元素，如果不将不可公度性内建为操作语义，如果不消除拷贝操作的不可能性，如果不将对偶性保留为结构性属性，不足以解决P vs.\ NP问题。}
	
	BSS模型将计算推向连续域，获得了强大的代数表达能力，但以偏离丘奇–图灵论题的离散预设为代价，且未解决操作不自足问题和对偶性缺失问题。CBTM/IVM框架则将计算保持在离散域内，与经典图灵机在计算能力上等价，但通过内建不可公度性、消除拷贝操作、保留对偶性丰富了操作语义。前者是``向外扩展''，后者是``向内深化''。P vs.\ NP问题的解决需要的不是更强大的计算能力，而是更丰富的语义表达能力——能够正面定义非确定性选择独立性、消除不可能操作、并保留其代数对偶性的概念资源。这正是CBTM/IVM框架所提供的东西。
	
	\section{``数计一体''作为方法论原则}
	
	前述诊断揭示了一个根本性的方法论困境：经典计算模型因其语法中立性而无法内在地表达非确定性计算所需的核心数学属性——不可公度性与对偶性，且经典NP定义要求了一个在该模型操作语义中不可实现的操作。本节提出一个规范性的方法论原则——数计一体——作为对这一困境的回应，并阐明CBTM/IVM框架如何通过参数化设计满足这一原则，从而系统性地修正第4-5节所诊断的操作不自足和三重概念错误。
	
	\subsection{原则的陈述}
	
	\textbf{数计一体}（the unity of mathematics and computation）主张：计算理论应当仿照数学，要求任意算法都应当具有明确的语义，并且语义与语法保持一致。
	
	在数学中，一个表达式不仅有其书写形式（语法），更有其指称的数学对象和满足的代数关系（语义）。数学推理的每一步都同时遵守语法规则和语义约束，两者始终协调一致。数学计算的正确性依赖于这种严格对应。
	
	经典计算理论背离了这一原则。图灵机模型刻意剥离了符号的语义，将计算还原为纯粹的语法操作。这种设计赋予了图灵机通用性，但也使计算模型丧失了``理解''其所处理对象的能力。经典NP定义不得不将非确定性的本质外包给外部量词和外部谓词，并在外包过程中扭曲了见证的本体论地位（函数指针被转换为数据，同时丢失了不可公度性）、必然消除了对偶性（拷贝不可能性的结构性后果）、且整个定义依赖于一个不可实现的操作（虚部到实部拷贝）——这正是这一背离的直接后果。
	
	数计一体原则要求重新统一语法与语义。具体而言，计算模型的设计应当满足以下要求：
	
	\begin{enumerate}
		\item \textbf{符号自带语义}：纸带上的每一个符号本身就是一个数学对象，而非等待外部赋予含义的裸标记。
		\item \textbf{操作保持语义}：每一步计算操作不仅移动符号，同时也在执行相应的数学运算，且运算结果忠实地反映在符号的语义变化中。
		\item \textbf{语义决定行为}：计算的走向（如是否分支）由符号的语义内容（如是否涉及不可公度元）决定，而非由外部量词或外部谓词事后赋予。
		\item \textbf{消除不可能操作}：计算模型的操作语义必须是自足的——模型不应要求其自身无法执行的操作（如虚部到实部的拷贝）。
		\item \textbf{对偶性显式保留}：非确定性选择的对偶结构应当作为操作语义的内在属性被保留，而非被编码和量词消除。
	\end{enumerate}
	
	在进入CBTM/IVM框架的具体机制之前，我们可以从一个更直观的代数视角来理解这一原则的动机。将经典图灵机的纸带看成是有理数域的编码，那么图灵机上的所有计算，本质上就是上一步和下一步之间进行仿射变换——每一步操作可以抽象为$X \leftarrow X + C$，其中$X$是当前状态和纸带内容的代数表示，$C$是由转移函数决定的常数。经典图灵机的全部表达能力，恰好被限制在有理数域内的仿射变换之中。
	
	从这个视角看，非确定性图灵机的非确定性操作可以自然地看作代数扩张。当非确定性图灵机在某个格局下面对多个后继选择时，这等价于在当前的基域中引入一个与已有元素不可公度的新生成元——正如从$\mathbb{Q}$扩张到$\mathbb{Q}(\sqrt{2})$时，$\sqrt{2}$无法用$\mathbb{Q}$中的任何有理线性组合来表示，因此机器必须``分叉''以同时探索包含和不包含该生成元的世界。
	
	由此，P是否等于NP的问题被转化为一个精确的代数问题：\textbf{代数数域上的计算是否可以由有理数域在多项式时间内模拟？}如果答案是肯定的，那么非确定性计算的代数扩张就是冗余的，可以在基域内被确定性模拟消除，即$\mathbf{P} = \mathbf{NP}$。如果答案是否定的，那么代数扩张提供了本质的计算资源，无法被基域内的仿射变换所模拟，即$\mathbf{P} \neq \mathbf{NP}$。
	
	\subsection{CBTM/IVM如何满足这一原则：对操作不自足与三重错误的系统性修正}
	
	CBTM是数计一体原则的第一个实例，它的设计体现了参数化方法的精髓，并且可以精确地理解为对第4节诊断的操作不自足和第5节诊断的三重概念错误的系统性修正。
	
	\textbf{修正操作不自足：从根本上消除拷贝操作。} 在经典NP定义中，见证$c$从虚部被拷贝到实部，这一操作在经典操作语义中不可实现，导致经典NP定义在操作上不自足。在CBTM中，这一问题被从根本上消除：\textbf{函数指针的地址不再需要拷贝到实部，而是常驻于虚部}。非确定性分支由机器内部的分支触发公理直接执行——当虚部$b=1$时，机器强制分叉，激活/不激活的选择直接发生在虚部，无需外部观察者介入，无需从控制信息到数据的跨层次转换。\textbf{从虚部到实部的拷贝这一不可能操作被彻底消除，经典NP定义的操作不自足性得到根本修正。}
	
	\textbf{修正第一重错误：函数指针常驻虚部，不可公度性得以保留。} 虚部符号的语义身份是明确的——它是控制信息，标记``此处涉及不可公度元''，而非数据。实部（数据层）与虚部（控制层）的严格分离，恢复了图灵机架构中代码与数据分离的基本原则。见证不再是被拷贝到实部的数据，而是虚部上激活/不激活的布尔标记序列。更为重要的是，虚部中的不同生成元（如$\sqrt{2}$和$\sqrt{3}$）在基域上不可公度，这一代数属性\textbf{被操作语义内在地保留}——CBTM通过分支触发公理强制分支的两条路径分别绑定于在基域上不可公度且互为对偶的生成元$\alpha$和$\beta$，从而保证了非确定性选择的独立性。这与经典NP定义形成鲜明对比：后者在将见证拷贝到实部的过程中，不可公度性已被完全抹除。
	
	\textbf{修正第二重错误：对偶性显式保留。} 在经典NP定义中，对偶性因编码而必然丢失。在CBTM/IVM中，\textbf{对偶性被显式保留在虚部空间中}。在CBTM中，非确定性分支的两条路径分别绑定于$\alpha$和$\beta$，这两个生成元在基域上不可公度，且由伽罗瓦自同构$\sigma: \alpha \leftrightarrow \beta$精确地相互映射。在IVM中，每次分支分配独立的生成元$\sqrt{p_i}$，其对偶为$-\sqrt{p_i}$。对偶性不再是被编码过程必然抹去的隐性结构，而是被内建为操作语义的显式代数事实。
	
	\textbf{修正第三重错误：语义内建而非外包。} 在经典NP定义中，非确定性由外部量词``$\exists c$''来断言，且依赖于不可实现的拷贝操作。在CBTM/IVM中，\textbf{非确定性由机器内部的代数机制来定义和执行}：分支触发公理规定当虚部$b=1$时机器强制分叉——这不是``存在多个后继''的语法约定，而是由不可公度性触发的代数行为；分支-激活引理规定每次分支强制激活一个独立的生成元，维度消耗是操作语义的强制性记账，无法被绕过。
	
	\textbf{参数化设计的精妙之处。} CBTM的参数化设计基于一个极为简洁的观察：经典非确定性图灵机（NTM）和确定性图灵机（DTM）的代码结构本质上是相同的，唯一的区别在于NTM在某些格局下允许转移关系中有多个五元组。CBTM对这一结构做了一个巧妙的一一映射，\textbf{完全没有改变TM或NTM的任何代码}：
	
	\begin{itemize}
		\item 对于DTM的每一个确定性步骤，CBTM将其映射为虚部$b=0$的符号操作——纸带上的内容被解释为基域中的元素，计算在基域内进行，不引入代数扩张。
		\item 对于NTM的每一个非确定性分叉步骤，CBTM将其映射为虚部$b=1$的符号操作——在分叉发生之前，当前纸带格子上的虚部被\textbf{预先标记}为$1$，表示该位置涉及域扩张中的不可公度元。两条分支路径分别绑定于对偶生成元。\textbf{分支直接由虚部标记触发，无需将控制信息拷贝到实部。}
	\end{itemize}
	
	这一映射是双射：给定任何一台DTM或NTM，存在唯一对应的CBTM（代码相同，仅在分叉处加标记）；反之，给定任何一台CBTM，抹去所有虚部标记即可恢复原DTM或NTM。\textbf{代码本质相同，虚部仅作为分支标记。}这正是参数化等价定理在操作层面的对应物——$\mathbf{P}_{cb} = \mathbf{P}$，$\mathbf{NP}_{cb} = \mathbf{NP}$。
	
	\textbf{参数化等价定理的双重意义。} 该定理不仅确立了新框架与经典框架的相关性（外延等价），更重要的是揭示了\textbf{内涵增强}的可能性。经典NTM对CBTM的模拟保持了计算能力，但模拟过程中必须将虚部所携带的代数语义（不可公度性、对偶性）\textbf{编码}为状态信息或额外的纸带符号。CBTM中由代数结构直接触发的分支，在经典模拟中变成了纯粹的语法操作——不同分支之间的代数独立性在经典模型的语法层面完全不可见。经典模拟器必须通过额外的计算开销来弥补它所无法表达且无法执行的操作（虚部直接触发分支、不可公度性保留、对偶性保留）。因此，参数化定理确立了两件事：新框架与经典框架识别相同的语言类（相关性），以及新框架拥有经典框架所不具备的语义分析工具和操作能力（优越性）。
	
	正是这一参数化设计，使得非确定性选择的独立性和对偶性获得了正面的代数刻画。两个分支之所以是``真正不同的选择''，不是因为它们对应不同的五元组（语法层面），而是因为它们分别绑定了对偶的不可公度生成元（代数层面）。这种独立性由域扩张的代数结构严格保证，对偶性由伽罗瓦群的作用精确刻画，无法通过基域内的任何运算相互转化。
	
	IVM将这一原则进一步推广到动态维度追踪。在IVM中，每次非确定性分支触发时，系统自动分配一个新的素数平方根地址$\sqrt{p_i}$，并在激活路径上记录值$1$。这一记账过程是IVM的定义性行为，独立于底层转移函数的设计。维度消耗因此成为语义记账的强制性结果，而非验证器的主动选择。
	
	\subsection{与丘奇-图灵论题的关系}
	
	数计一体可以被视为丘奇-图灵论题的一个自然延伸。丘奇-图灵论题断言任何算法可计算的函数都是图灵机可计算的。数计一体则进一步要求：任何数学对象所携带的语义结构（如不可公度性、线性无关性、代数对偶性等）都应当在计算模型的操作语义中获得正面的、内在的表达，且计算模型的操作语义必须是自足的——不应依赖其自身无法执行的操作。
	
	这不是对丘奇-图灵论题的否定，而是对其``语法中立性''预设的修正。丘奇-图灵论题关注的是``能计算什么''的外延等价；数计一体关注的是``如何计算''的内涵充分性和操作自足性。一个计算模型不应仅满足于可计算性的外延等价，而应当追求语义表达的内蕴充分性和操作自足性。
	
	\subsection{范式转换：从语法范式到语义范式}
	
	本文所倡导的范式转换——从``语义外包、操作不自足''到``语义内建、操作自足''——或许不仅适用于P vs.\ NP问题。其他复杂性类的分离问题、随机性与非确定性的关系问题、以及量子计算中的类似问题，都可能从这一方法论转换中获益。这些方向留待未来研究。
	
	\section{结论}
	
	\subsection{诊断的总结}
	
	本文提出了一个关于P vs.\ NP问题的哲学诊断。其核心论证可以概括为以下链条：
	
	\begin{enumerate}
		\item 经典NP定义将非确定性的本质外包给外部量词``存在见证''（$\exists c$）和外部谓词$IsWitness(c)$（双重外包）。
		\item 更为根本的是，经典NP定义隐式地要求了一个在经典操作语义中\textbf{不可实现}的操作——将虚部控制信息拷贝到实部成为数据$c$。这使得经典NP定义是\textbf{操作不自足}的：它诉诸了一个其自身形式体系无法提供的资源。在实际的经典NP验证中，$c$由外部观察者直接提供，``非确定性''因此不是机器内部的操作特征，而是外部语义标签。
		\item 在这一拷贝不可能性的基础上，经典NP定义必然导致三重概念错误：(a) 见证$c$从虚部控制信息被扭曲为实部数据，同时导致\textbf{不可公度性的丢失}——比特间代数独立性被抹除（第一重错误）；(b) 对偶性被存在量词消除——从虚部到实部的编码必然丢失对偶结构（第二重错误）；(c) 因为不可公度性（第一重错误的后果）和对偶性（第二重错误的后果）都已丢失，经典模型无法正面刻画非确定性选择的独立性和结构，只能将其外包给外部量词和外部谓词，导致非确定性概念的\textbf{语义残缺}（第三重错误）。
		\item 外部谓词$IsWitness(c)$的内涵需要不可公度性来填充——见证的每一个比特代表一次独立的函数指针选择，独立性由不可公度性保证，结构由对偶性刻画。
		\item 经典图灵机的操作语义无法内在地表达不可公度性（不可公度性元定理）。经典模型可以在弱意义上判定不可公度性，但无法在强意义上将其锚定。非平凡编码是外延指代而非内涵内建。
		\item 因此，经典NP定义在语义上是不完整的，在操作上是不自足的：它指明了哪些问题是NP的（外延），但没有正面刻画非确定性的本质是什么（内涵），在定义过程中必然丢失了非确定性的核心语义属性（不可公度性和对偶性），且整个定义依赖于一个不可实现的操作。
		\item 聪明算法假设（$\mathbf{P}=\mathbf{NP}$）要求用无法表达不可公度性、无法执行虚部到实部拷贝、且丧失了对偶性的内部语言，去等价地实现依赖于不可公度性、依赖于该拷贝结果、且依赖于对偶性的外部语义。这是一个语义层面和操作层面的双重资源错位。
		\item 这解释了为什么三大障碍在经典框架内长期难以逾越。
	\end{enumerate}
	
	\subsection{从消解到解决}
	
	基于上述诊断，经典P vs.\ NP问题在经典框架内难以获得充分的表达资源——它要求用缺乏必要概念资源、存在操作不自足、且存在概念扭曲的语言去讨论需要这些概念的问题。这不是说该问题没有真值，而是说其证明需要超越经典模型的概念框架。
	
	CBTM/IVM框架提供了这样的概念框架。通过将不可公度性内建为操作语义的内在属性、让对偶性常驻虚部、从根本上消除拷贝操作的必要性、以语义内建取代语义外包，该框架系统性地修正了经典NP定义的操作不自足和三重概念错误。在这一新框架中，P与NP都获得了内在的代数定义：P类对应零维验证器（$b=0$，无需代数扩张），NP类对应非零维验证器（$b=1$，涉及代数扩张）。两者的差异被转化为一个精确的代数问题——域扩张是否为非确定性计算提供了本质的代数资源。这一思路类似于非欧几何对平行公设问题的解决：不是用旧方法证明旧命题，而是构建了能够重新表述问题的新框架。
	
	\subsection{对计算基础的意义}
	
	本文的诊断和解决路径表明，计算理论的基础可能需要在``语义内建''而非``语法中立''的方向上重新审视。经典图灵机模型的语法中立性是其力量的源泉，也是其表达能力的边界所在。当我们面对需要正面定义``独立性''或``不可归约性''的复杂性理论问题时，纯粹语法层面的描述可能是不够的。
	
	更为根本的是，本文揭示了一个被长期忽视的事实：经典NP定义不仅缺失了必要的语义资源，而且\textbf{在操作语义上是不自足的}——它要求了一个其自身形式体系无法执行的操作。这一操作不自足性是经典NP定义所有结构性缺陷的最终根源。由此必然导致三重概念错误：第一重错误（函数指针→数据）丢失了不可公度性，第二重错误（编码与量词）丢失了对偶性，第三重错误（语义外包）是前两重在语义层面的必然后果。这一诊断不仅解释了为什么经典框架无法完成P vs NP的证明，也为未来计算模型的设计提供了方法论指导：一个完备的计算模型，必须遵守``数计一体''原则，保持语法与语义的统一，确保操作语义的自足性，让不可公度性内建，让对偶性常驻，从根本上消除对不可能操作的依赖。
	
	本文所倡导的范式转换——从语义外包到语义内建，从操作不自足到操作自足——或许不仅适用于P vs.\ NP问题。其他复杂性类的分离问题、随机性与非确定性的关系问题、以及量子计算中的类似问题，都可能从这一方法论转换中获益。这些方向留待未来研究。
	
	\hfill $\square$
	
	\bibliographystyle{plain}
	\bibliography{../cb}
	
\end{document}